\clearpage
\fancyhead[L]{\small COMET.}
\SourceLeaf{447}
\begin{center}
{\large COMET.}

\noindent\rule{\linewidth}{0.4pt}

Göttingische gelehrte Anzeigen. Stück 26. Seite 249..251. 1823 Febr. 15.

\noindent\rule{\linewidth}{0.4pt}
\end{center}

Ein Schreiben des Herrn \textsc{Rümker} aus Paramatta in New South Wales an Herrn Hofr. \textsc{Gauss} vom
August 1822 enthält ausser vielen andern daselbst angestellten astronomischen Beobachtungen,
deren Bekanntmachung für einen andern Ort aufgespart wird, auch die Beobachtungen des
\textsc{Encke}schen Cometen, welche Hr. \textsc{Rümker} im Juni v.\,J. angestellt hat. Bekanntlich war die
Hoffnung, dass dieser Comet bei seiner im Jahr 1822 erfolgten Wiederkehr zur Sonne in Europa
observirt werden könnte, sehr schwach, und in der That sind auch alle Bemühungen europäischer
Astronomen, ihn zu sehen, ganz fruchtlos gewesen. Dagegen rechnete man darauf, dass er den
Beobachtern in der südlichen Hemisphäre nicht entgehen würde: allein, da man bereits Nachricht
hat, dass die Nachforschungen der englischen Astronomen auf dem Vorgebirge der guten Hoffnung
\emph{nicht gelungen} sind, so sah man um so sehnlicher den Nachrichten aus Neuholland entgegen,
und wir freuen uns daher um so mehr, in unsern Blättern, die im 83. Stück von 1819 [S.~420 d.~B.]
die erste Nachricht von Hrn. \textsc{Encke}'s so höchst merkwürdiger Entdeckung gegeben haben, nun die
durch unsern geschickten Landsmann gewonnene Bestätigung derselben mittheilen zu können.
Folgendes sind die \emph{sämmtlichen} Beobachtungen, die Herr \textsc{Rümker} gemacht hat, da nach dem
23. Juni der Mondschein hinderlich, und nach dem Vollmonde der Comet zu lichtschwach war, um
noch ferner beobachtet werden zu können.

\Needspace{23\baselineskip}
\begin{center}
Beobachtungen des \textsc{Encke}schen Cometen in \emph{Paramatta}.
\end{center}

\begin{center}
\small\setlength{\tabcolsep}{3pt}
\begin{tabular}{l|ccc|ccc|ccc}
\multicolumn{1}{c|}{1822} & \multicolumn{3}{c|}{Sternzeit} & \multicolumn{3}{c|}{Mittlere AR.} & \multicolumn{3}{c}{Mittlere Declination} \\
\hline
Juni 2 & $10^{\mathrm{u}}$ & $39^{\mathrm{m}}$ & $25^{\mathrm{s}}$ & $92^{\circ}$ & $43'$ & $51''3$ & $17^{\circ}$ & $39'$ & $46''3$ Nordl. \\
3 & 11 & --- & --- & 93 & 46 & 20.7 & 16 & 53 & 7.5 \\
4 & 11 & 3 & 0 & 94 & 46 & 0.0 & 16 & 4 & 36.7 \\
6 & 11 & 7 & 38 & 96 & 42 & 11.6 & 14 & 22 & 42.0 \\
7 & 11 & 3 & 10 & 97 & 38 & 15 & 13 & 26 & 5 \\
8 & 11 & 17 & 25 & 98 & 33 & 47.7 & 12 & 31 & 18.6 \\
10 & 11 & 20 & 0 & 100 & 24 & 43.8 & 10 & 29 & 49.5 \\
11 & 11 & 24 & 39 & 101 & 19 & 44.5 & 9 & 26 & 4.6 \\
12 & 11 & 40 & 0 & 102 & 17 & 52 & 8 & 18 & 30 \\
13 & 11 & 42 & 4 & 103 & 15 & 2 & 7 & 6 & 30 \\
14 & 11 & 55 & 0 & 104 & 15 & 40 & 5 & 52 & 27 \\
15 & 11 & 40 & 48 & 105 & 17 & 0.5 & 4 & 33 & 40 \\
19 & 12 & 13 & 38 & 109 & 54 & 36.4 & 1 & 29 & 43.7 Südl. \\
20 & 12 & 16 & 53 & 111 & 14 & 26.9 & 3 & 14 & 29.1 \\
22 & 13 & 18 & 46 & 114 & 12 & 20.5 & 7 & 8 & --- \\
23 & 12 & 53 & 55 & 115 & 47 & 41.7 & 9 & 9 & 48.4 \\
\end{tabular}
\end{center}

Ob der Zusatz \emph{mittlere} Rectascension und Declination auf ein Mittel mehrerer
Beobachtungen, oder auf eine Befreiung von der Nutation Bezug haben soll, wird nicht bemerkt.
Zur Vergleichung sind Sterne aus \textsc{Piazzi}'s Catalog und aus der \emph{Histoire Céleste} gebraucht,
worüber Hr. \textsc{Rümker} das nähere mit erster Gelegenheit nachzuliefern verspricht.

Die Polhöhe der Sternwarte in Paramatta findet Herr \textsc{Rümker} aus den beiden Solstitien von
December 1821 und von Juni 1822

\[
33^{\circ}\,48'\,41''97
\]

\Needspace{16\baselineskip}
Aus einer Bedeckung von $\alpha$ Scorpii, am 10. April 1822, wo

% 
\SourceLeaf{448}
\begin{center}
\small
\begin{tabular}{l l}
der Eintritt \mbox{.\,.\,.\,.\,.\,.\,.\,.\,.\,.\,.} & $18^{\mathrm{u}}\;35^{\mathrm{m}}\;47^{\mathrm{s}}4$ Mittlere Zeit \\
der Austritt \mbox{.\,.\,.\,.\,.\,.\,.\,.\,.\,.\,.\,.} & $19\quad14\quad27.9$
\end{tabular}
\end{center}

\noindent beobachtet wurde, findet Hr. \textsc{Rümker}, die Mondsörter aus dem Nautical Almanac entlehnend, die Länge von Paris
\begin{center}
\begin{tabular}{l l}
10 St. $3^{\mathrm{m}}\;56^{\mathrm{s}}3$ & aus dem Eintritt \\
$10\quad4\quad7.3$ & aus dem Austritt.
\end{tabular}
\end{center}
